\section{Related Work}
\label{section:2}

\acresetall % Reset acronyms so they are printed in full in each new section.

Multiple visualization-adjacent literatures have studied methods for describing charts and graphics through natural language\,---\,including accessible media research, \ac{HCI}, \ac{CV}, and \ac{NLP}. But, these various efforts have been largely siloed from one another, adopting divergent methods and terminologies (e.g., the terms ``caption'' and ``description'' are used inconsistently).
% \footnote{
    % The terms ``caption'' and ``description'' are used inconsistently across these literatures.
    % ``caption'' is often used by researchers working with computer vision methods~\cite{qian_formative_2020, xu_show_2016}, whereas ``description'' is more common in accessible media research and \ac{HCI}~\cite{gould_effective_2008, morash_guiding_2015}.
% }
%     %as well consumer-grade visualization software, such as Tableau
%     % In the latter literatures, ``description'' is preferred to ``caption'' because a visualization description may \emph{contain} a caption.
%     % For example, a visualization may come paired with a natural language caption.
%     % But, when providing a natural language \emph{alternative} to that visualization (specifically, for people with visual disabilities), the caption
%     % % paired with the visualization
%     % may itself be \emph{included in} the description. Thus, ``description'' is conceptualized more broadly than ``caption''.
%     In this paper we use ``captioning'' to refer to automatic methods, and ``describing'' to refer to the more general practice of conveying visualized content through natural language.
%     % Additionally, following Bertin, we use ``content'' to refer to the {concepts and information} that can be conveyed through either natural language, visualization, or both.
% }
Here, we survey the diverse terrain of literatures intersecting visualization and natural language.

%----------------------------------------------------------------
\subsection{Automatic Methods for Visualization Captioning}
%----------------------------------------------------------------

Automatic methods for generating visualization captions broadly fall into two categories: those using \ac{CV} and \ac{NLP} methods when the chart is a rasterized image (e.g., \textsc{jpeg}s or \textsc{png}s); and those using structured specifications of the chart's construction (e.g., grammars of graphics).

\subsubsection{Computer Vision and Natural Language Processing}
Analogous to the long-standing \ac{CV} and \ac{NLP} problem of automatically captioning photographic images~\cite{lin_microsoft_2014, krishna_visual_2017, karpathy_deep_2017}, recent work on visualization captioning has aimed to automatically generate accurate and descriptive natural language sentences for charts~\cite{chen_neural_2019, chen_figure_2019, chen_figure_2020, balaji_chart-text_2018, obeid_chart--text_2020, lai_automatic_2020, qian_generating_2021}.
Following the encoder-decoder framework of statistical machine translation~\cite{vinyals_show_2015, xu_show_2016}, these approaches usually take rasterized images of visualizations as input to a \ac{CV} model (the encoder), which learns the visually salient features for outputting a relevant caption via a language model (the decoder).
Training data consists of \textlangle chart, caption\textrangle\, pairs, collected via web-scraping and crowdsourcing~\cite{qian_formative_2020}, or created synthetically from pre-defined sentence templates~\cite{kahou_figureqa_2018}.
While these approaches are well-intentioned, in aiming to address the engineering problem of \emph{how} to automatically generate natural language captions for charts, they have largely sidestepped the complementary (and prior) question: \emph{which} semantic content should be generated to begin with?
Some captions may be more or less descriptive than others, and different readers may receive different semantic content as more or less useful, depending on their levels of data literacy, domain-expertise, and/or visual perceptual ability~\cite{macleod_understanding_2017, morash_guiding_2015, moraes_evaluating_2014}.
To help orient work on automatic visualization captioning, our four-level model of semantic content offers a means of asking and answering these more human-centric questions.

\subsubsection{Structured Visualization Specifications}
In contrast to rasterized images of visualizations, chart templates~\cite{team_bokeh_2014}, component-based architectures~\cite{geveci_vtk_2012}, and grammars of graphics~\cite{satyanarayan_vega-lite_2017} provide not only a structured representation of the visualization's construction, but typically render the visualization in a structured manner as well.
For instance, most of these approaches either render the output visualization as \ac{SVG} or provide a scenegraph API.
Unfortunately, these output representations lose many of the semantics of the structured input (e.g., which elements correspond to axes and legends, or how nesting corresponds to visual perception).
As a result, most present-day visualizations are inaccessible to people who navigate the web using screen readers.
For example, using Apple's VoiceOver to read D3 charts rendered as \ac{SVG} usually outputs an inscrutable mess of screen coordinates and shape rendering properties.
Visualization toolkits can ameliorate this by leveraging their structured input to automatically add \ac{ARIA} attributes to appropriate output elements, in compliance with the \ac{W3C}'s \ac{WAI} guidelines~\cite{w3c_wai_2019}.
Moreover, this structured input representation can also simplify automatically generating natural language captions through template-based mechanisms, as we discuss in \S~\ref{sec:level1}.

%----------------------------------------------------------------
\subsection{Accessible Media and Human-Computer Interaction}
%----------------------------------------------------------------

While automatic methods researchers often note accessibility as a worthy motivation~\cite{demir_summarizing_2012, obeid_chart--text_2020, qian_formative_2020, qian_generating_2021, elzer_exploring_2005, elzer_browser_2007}, evidently few have collaborated directly with disabled people~\cite{choi_visualizing_2019, moraes_evaluating_2014} or consulted existing accessibility guidelines~\cite{lundgard_sociotechnical_2019}.
Doing so is more common to \ac{HCI} and accessible media literatures~\cite{morris_rich_2018, sharif_understanding_2021}, which broadly separate into two categories corresponding to the relative expertise of the description authors: those authored by {experts} (e.g., publishers of accessible media) and those authored by {non-experts} (e.g., via crowdsourcing or online platforms).

\subsubsection{Descriptions Authored by Experts}
Publishers have developed guidelines for describing graphics appearing in \ac{STEM} materials~\cite{gould_effective_2008, benetech_making_nodate}.
Developed by and for authors with some expert accessibility knowledge, these guidelines provide best practices for conveying visualized content in traditional media (e.g., printed textbooks, audio books, and tactile graphics).
But, many of their prescriptions\,---\,particularly those relating to the \emph{content} conveyed by a chart, rather than the \emph{modality} through which the chart is rendered\,---\,are also applicable to web-based visualizations.
Additionally, web accessibility guidelines from \ac{W3C} provide best-practices for writing descriptions of ``complex images'' (including canonical chart types), either in a short description alt text attribute, or as a long textual description displayed alongside the visual image~\cite{w3c_wai_2019}.
While some of these guidelines have been adopted by visualization practitioners~\cite{cesal_writing_2020, schepers_why_2020, elavsky_chartability_2021, fisher_creating_2019, watson_accessible_2018, watson_accessible_2017-1, fossheim_introduction_2020}, we here bring special attention to the empirically-grounded and well-documented guidelines created by the \textsc{wgbh} National Center for Accessible Media~\cite{gould_effective_2008} and by the Benetech Diagram Center~\cite{benetech_making_nodate}.

\subsubsection{Descriptions Authored by Non-Experts}
Frequently employed in \ac{HCI} and visualization research, crowdsourcing is a technique whereby remote non-experts complete tasks currently infeasible for automatic methods, with applications to online accessibility~\cite{bigham_vizwiz_2010}, as well as remote description services like \emph{Be My Eyes}.
For example, Morash et al. explored the efficacy of two types of non-expert tasks for authoring descriptions of visualizations: non-experts authoring free-form descriptions without expert guidance, versus those filling-in sentence templates pre-authored by experts~\cite{morash_guiding_2015}.
While these approaches can yield more richly detailed and ``natural''-sounding descriptions (as we discuss in \S~\ref{section:5}), and also provide training data for auto-generated captions and annotations~\cite{qian_formative_2020, kong_extracting_2014}, it is important to be attentive to potential biases in human-authored descriptions~\cite{bennett_its_2021}.

%----------------------------------------------------------------
\subsection{Natural Language Hierarchies and Interfaces}
%----------------------------------------------------------------

Apart from the above methods for generating descriptions, prior work has adopted linguistics-inspired framings to elucidate how natural language is used to {describe}\,---\,as well as {interact with}\,---\,visualizations.

\subsubsection{Using Natural Language to Describe Visualizations}

Demir et al. have proposed a hierarchy of six syntactic complexity levels corresponding to a set of propositions that might be conveyed by bar charts~\cite{demir_summarizing_2012}.
Our model differs in that it orders \emph{semantic} content\,---\,i.e., \emph{what} meaning the natural language sentence conveys\,---\,rather than \emph{how} it does so syntactically. 
Thus, our model is agnostic to a sentence’s length, whether it contains multiple clauses or conjunctions, which has also been a focus of prior work in automatic captioning~\cite{qian_formative_2020}.
Moreover, whereas Demir et al. speculatively ``envision'' their set of propositions to construct their hierarchy, we arrive to our model empirically through a multi-stage grounded theory process~\sectionlabel{section:3}.
Perhaps closest to our contribution are a pair of papers by Kosslyn~\cite{kosslyn_understanding_1989} and Livingston \& Brock~\cite{livingston_position_2020}.
Kosslyn draws on canonical linguistic theory, to introduce three levels for analyzing charts: the \emph{syntactic} relationship between a visualization elements; the \emph{semantic} meaning of these elements in what they depict or convey; and the \emph{pragmatic} aspects of what these elements convey in the broader context of their reading~\cite{kosslyn_understanding_1989}.
We seeded our model construction with a similar linguistics-inspired framing, but also evaluated it empirically, to further decompose the semantic levels~\sectionlabel{section:3.1}.
Livingston \& Brock adapt Kosslyn's ideas to generate what they call ``visual sentences'': natural language sentences that are the result of executing a single, specific analytic task against a visualization~\cite{livingston_position_2020}.
Inspired by the Sentence Verification Technique (\textsc{svt})~\cite{royer_sentence_1979, royer_developing_2001}, this work considers visual sentences for assessing graph comprehension, hoping to offer a more ``objective'' and automated alternative to existing visualization literacy assessments~\cite{lee_vlat_2016, galesic_graph_2011}.
While we adopt a more qualitative process for constructing our model, Livingston \& Brock's approach suggests opportunities for future work: might our model map to similarly-hierarchical models of analytic tasks~\cite{brehmer_multi-level_2013, amar_low-level_2005}?

\subsubsection{Using Natural Language to Interact with Visualizations}

Adjacently, there is a breadth of work on \ac{NLIs} for constructing and exploring visualizations~\cite{narechania_nl4dv_2021, hearst_toward_2019, setlur_inferencing_2019, kim_answering_2020}.
While our model primarily considers the natural language sentences that are \emph{conveyed by} visualizations (cf., natural language as \emph{input} for chart specification and exploration)~\cite{srinivasan_collecting_2021}, our work may yet have implications for \ac{NLIs}.
For example, Hearst et al. have found that many users of chatbots prefer \emph{not} to see charts and graphics alongside text in the conversational dialogue interface~\cite{hearst_would_2019}.
By helping to decouple visual-versus-linguistic data representations, our model might be applied to offer these users a textual alternative to inline charts.
Thus, we view our work as complementary to \ac{NLIs}, facilitating multimodal and more accessible data representations~\cite{kim_facilitating_2018}, while helping to clarify the theoretical relationship between charts and captions~\cite{kim_towards_2021, ottley_curious_2019}, and other accompanying text~\cite{xiong_curse_2020, kong_frames_2018, kong_trust_2019, adar_communicative_2020}.